% ============================================================
% Lecture 8 -- Fundamentals of Linear Models and the estimation problem
% ============================================================
\section{Fundamentals of Linear Models and the Estimation Problem}\label{sec:L08}

\lectureheader{8}{9RbRhz776-k}{Week-3 handwritten notes \texttt{Feb9.pdf}, first part (``Model 1'', within/between variation, design matrix, identifiability)}

\begin{guiding*}
By the end of this lecture you should be able to:
\begin{itemize}
  \item write down the one-way classification model $Y_{ij}=\mu+\mu_i+\eps_{ij}$ and state the three estimation questions for it;
  \item define $\SSW$, $\SSB$ and $\TSS$, prove $\TSS=\SSW+\SSB$, and state and prove their distributions and independence;
  \item write the model as $Y=X\beta+\eps$, find $\rank(X)$, and explain why its parameters are \emph{not identifiable};
  \item show that contrasts such as $\mu_1-\mu_2$ are identifiable, and that a constraint like $\sum\mu_i=0$ makes all parameters identifiable.
\end{itemize}
\end{guiding*}

\subsection{Motivation}
Week~2 fitted lines to data. From now on we develop the \emph{theory}: when can parameters be
estimated, what is the best estimate, and how are its errors distributed? The instructor begins
with the model that underlies the analysis of variance.

\subsection{The one-way classification model}
\begin{definition}[Model 1: one-way classification]\label{def:L08-oneway}
There are $k$ populations (groups). From population $i$ we observe $n_i$ values $Y_{i1},\dots,Y_{in_i}$:
\[
Y_{ij}=\mu+\mu_i+\eps_{ij},\qquad 1\le i\le k,\quad 1\le j\le n_i,
\]
where $\mu,\mu_1,\dots,\mu_k\in\RR$ are unknown and the errors $\eps_{ij}\sim\N(0,\sigma^2)$ are
independent. The total sample size is $n=n_1+n_2+\dots+n_k$. Here $\mu$ is a general
level and $\mu_i$ the \vocab{effect} of group $i$. The mean of population $i$ is $\mu+\mu_i$.
\end{definition}
The data can be pictured as $k$ columns:
\begin{center}
\begin{tabular}{cccc}\hline
Population 1 & Population 2 & $\cdots$ & Population $k$\\\hline
$Y_{11}$ & $Y_{21}$ & & $Y_{k1}$\\
$\vdots$ & $\vdots$ & & $\vdots$\\
$Y_{1n_1}$ & $Y_{2n_2}$ & & $Y_{kn_k}$\\\hline
\end{tabular}
\end{center}

\begin{question}\label{q:L08-questions}
The lecture poses three questions.
\begin{enumerate}
  \item[(i)] Can we estimate $\mu$? If yes, what is the estimate?
  \item[(ii)] Can we estimate $\mu+\mu_i$? If yes, what is the estimate?
  \item[(iii)] Can we estimate $\mu_i-\mu_j$ for $1\le i\ne j\le k$? If yes, what is the estimate?
\end{enumerate}
\end{question}
The answers are (i) no, (ii) yes, (iii) yes. This lecture shows why (i) fails and gives
unbiased estimates for (iii). \Cref{sec:L10,sec:L13} give the general theory.

\subsection{Within-group and between-group variation}
Write $\bar Y_{i\cdot}=\frac1{n_i}\sum_{j=1}^{n_i}Y_{ij}$ for the mean of group $i$ and
$\bar Y_{\cdot\cdot}=\frac1n\sum_{i=1}^k\sum_{j=1}^{n_i}Y_{ij}$ for the grand mean.

\begin{definition}[Sums of squares]\label{def:L08-SS}
\begin{itemize}
  \item The variance of group $i$ is estimated unbiasedly by the $i$th sample variance
  $S_i^2=\frac{1}{n_i-1}\sum_{j=1}^{n_i}(Y_{ij}-\bar Y_{i\cdot})^2$ (for $n_i\ge 2$).
  \item The \vocab{sum of squares within} groups is $\SSW=\sum_{i=1}^k\sum_{j=1}^{n_i}(Y_{ij}-\bar Y_{i\cdot})^2$.
  \item The \vocab{sum of squares between} groups is $\SSB=\sum_{i=1}^k n_i(\bar Y_{i\cdot}-\bar Y_{\cdot\cdot})^2$.
  \item The \vocab{total sum of squares} is $\TSS=\sum_{i=1}^k\sum_{j=1}^{n_i}(Y_{ij}-\bar Y_{\cdot\cdot})^2$.
\end{itemize}
\end{definition}

\begin{theorem}[Decomposition of the total sum of squares]\label{thm:L08-decomp}
$\TSS=\SSW+\SSB$.
\end{theorem}
\begin{proof}
Write $Y_{ij}-\bar Y_{\cdot\cdot}=(Y_{ij}-\bar Y_{i\cdot})+(\bar Y_{i\cdot}-\bar Y_{\cdot\cdot})$ and square:
\[
\TSS=\sum_{i,j}(Y_{ij}-\bar Y_{i\cdot})^2+\sum_{i,j}(\bar Y_{i\cdot}-\bar Y_{\cdot\cdot})^2
+2\sum_{i=1}^k(\bar Y_{i\cdot}-\bar Y_{\cdot\cdot})\sum_{j=1}^{n_i}(Y_{ij}-\bar Y_{i\cdot}).
\]
The inner sum in the cross term is $\sum_j Y_{ij}-n_i\bar Y_{i\cdot}=0$ for every $i$, so the
cross term vanishes. The first sum is $\SSW$. In the second the summand does not depend on
$j$, so it equals $\sum_i n_i(\bar Y_{i\cdot}-\bar Y_{\cdot\cdot})^2=\SSB$.
\end{proof}

The notes then list distributional facts, the first two posed as exercises. To prove them we
need two tools about normal vectors.

\begin{lemma}[Quadratic forms in projections]\label{lem:L08-chisq}
Let $Z\sim\N(0,I_m)$ (i.e.\ $Z_1,\dots,Z_m$ i.i.d.\ $\N(0,1)$) and let $A$ be an $m\times m$
symmetric matrix with $A^2=A$ and $\rank(A)=r$. Then $Z\T AZ\sim\chi^2_r$.
\end{lemma}
\begin{proof}
By the spectral theorem, $A=Q\Lambda Q\T$ with $Q$ orthogonal and $\Lambda$ diagonal. From
$A^2=A$ we get $\Lambda^2=\Lambda$, so every eigenvalue is $0$ or $1$. The number of $1$'s is
$\rank(\Lambda)=\rank(A)=r$. Order them first. Put $W=Q\T Z$. $W$ is a linear function of a
normal vector, so it is normal, with mean $0$ and covariance $Q\T I Q=I$. Hence
$W_1,\dots,W_m$ are i.i.d.\ $\N(0,1)$, and
$Z\T AZ=W\T\Lambda W=W_1^2+\dots+W_r^2\sim\chi^2_r$.
\end{proof}

\begin{lemma}[Uncorrelated jointly normal vectors are independent]\label{lem:L08-indep}
If $(U,V)$ is jointly normal and $\Cov(U,V)=0$ (every component of $U$ is uncorrelated with
every component of $V$), then $U$ and $V$ are independent. Hence so are $g(U)$ and $h(V)$ for any
functions $g,h$.
\end{lemma}
\begin{proof}
Let $m_U,m_V$ be the means and $\Sigma_U,\Sigma_V$ the covariance matrices. The joint covariance
matrix is block diagonal, $\mathrm{diag}(\Sigma_U,\Sigma_V)$. The moment generating function of a
normal vector $W$ with mean $m$ and covariance $\Sigma$ is $\EE e^{t\T W}=\exp(t\T m+\tfrac12 t\T\Sigma t)$.
For $t=(s,u)$ the block structure gives
\[
\EE e^{s\T U+u\T V}=\exp\bigl(s\T m_U+\tfrac12 s\T\Sigma_Us\bigr)\exp\bigl(u\T m_V+\tfrac12 u\T\Sigma_Vu\bigr)
=\EE e^{s\T U}\,\EE e^{u\T V}.
\]
A joint moment generating function that factorises determines a product distribution, so $U\indep V$.
\end{proof}

\begin{theorem}[Distribution of the sums of squares]\label{thm:L08-dist}
In Model~1:
\begin{enumerate}
  \item $\SSW/\sigma^2\sim\chi^2_{n-k}$: $\SSW$ has $n-k$ \vocab{degrees of freedom};
  \item $\SSW$ is independent of $\SSB$;
  \item if $\mu_1=\mu_2=\dots=\mu_k$, then $\SSB/\sigma^2\sim\chi^2_{k-1}$ ($k-1$ degrees of freedom) and $\TSS/\sigma^2\sim\chi^2_{n-1}$ ($n-1$ degrees of freedom).
\end{enumerate}
\end{theorem}
\begin{proof}
(1) Fix $i$ and let $\eps_i=(\eps_{i1},\dots,\eps_{in_i})\T$. Since $Y_{ij}-\bar Y_{i\cdot}=\eps_{ij}-\bar\eps_{i\cdot}$,
\[
\sum_{j}(Y_{ij}-\bar Y_{i\cdot})^2=\eps_i\T\Bigl(I_{n_i}-\tfrac1{n_i}\one\one\T\Bigr)\eps_i .
\]
The matrix $C_i=I-\frac1{n_i}\one\one\T$ is symmetric and idempotent: $C_i^2=I-\frac2{n_i}\one\one\T+\frac1{n_i^2}\one(\one\T\one)\one\T=C_i$.
Its rank equals its trace, $n_i-1$. The trace equals the rank for idempotent matrices, because
the eigenvalues are $0$ or $1$. Applying \cref{lem:L08-chisq} to $Z=\eps_i/\sigma$ gives
$\sum_j(Y_{ij}-\bar Y_{i\cdot})^2/\sigma^2\sim\chi^2_{n_i-1}$. The $k$ groups are independent. A sum
of independent $\chi^2$ variables is $\chi^2$ with the degrees of freedom added: multiply the
moment generating functions $(1-2t)^{-(n_i-1)/2}$. So $\SSW/\sigma^2\sim\chi^2_{\sum(n_i-1)}=\chi^2_{n-k}$.

(2) $\SSW$ is a function of the vector $U=(\eps_{ij}-\bar\eps_{i\cdot})_{i,j}$. $\SSB$ is a function of
$V=(\bar Y_{1\cdot},\dots,\bar Y_{k\cdot})$, because $\bar Y_{\cdot\cdot}=\frac1n\sum n_i\bar Y_{i\cdot}$. Both are linear
functions of the normal vector $\eps$, so $(U,V)$ is jointly normal. For $i'\ne i$,
$\Cov(\bar Y_{i'\cdot},\eps_{ij}-\bar\eps_{i\cdot})=0$ because different groups are independent. For $i'=i$,
\[
\Cov(\bar\eps_{i\cdot},\eps_{ij}-\bar\eps_{i\cdot})=\Cov(\bar\eps_{i\cdot},\eps_{ij})-\Var(\bar\eps_{i\cdot})=\frac{\sigma^2}{n_i}-\frac{\sigma^2}{n_i}=0 .
\]
By \cref{lem:L08-indep}, $U\indep V$, hence $\SSW\indep\SSB$.

(3) Suppose all $\mu_i$ equal a common value, so every group has mean $\vartheta=\mu+\mu_1$. Then
$\bar Y_{i\cdot}\sim\N(\vartheta,\sigma^2/n_i)$ independently. Put $Z_i=\sqrt{n_i}(\bar Y_{i\cdot}-\vartheta)/\sigma$, so
$Z=(Z_1,\dots,Z_k)\T\sim\N(0,I_k)$. Let $u=(\sqrt{n_1},\dots,\sqrt{n_k})\T/\sqrt n$, a unit vector. By
\cref{prop:L01-bestconst} (with weights $n_i$, same proof),
$\sum_i n_i(\bar Y_{i\cdot}-\vartheta)^2=\SSB+n(\bar Y_{\cdot\cdot}-\vartheta)^2$. Dividing by $\sigma^2$:
\[
\|Z\|^2=\frac{\SSB}{\sigma^2}+(u\T Z)^2,\qquad\text{since } \sqrt n(\bar Y_{\cdot\cdot}-\vartheta)/\sigma=\sum_i\tfrac{\sqrt{n_i}}{\sqrt n}Z_i=u\T Z .
\]
So $\SSB/\sigma^2=Z\T(I_k-uu\T)Z$. The matrix $I_k-uu\T$ is symmetric and idempotent with trace
$k-1$, and \cref{lem:L08-chisq} gives $\chi^2_{k-1}$. Finally $\TSS=\SSW+\SSB$ is a sum of
independent $\chi^2_{n-k}$ and $\chi^2_{k-1}$ variables (times $\sigma^2$), hence $\TSS/\sigma^2\sim\chi^2_{n-1}$.
\end{proof}

\begin{note}
Degrees of freedom: $n-k$ for $\SSW$, $k-1$ for $\SSB$, $n-1$ for $\TSS$, and $(n-k)+(k-1)=n-1$.
Parts of the handwritten notes are hard to read at this point. The degrees of freedom stated
here are the standard ones, and they are the ones the proof gives.
\end{note}

\begin{example}[Sums of squares by hand]\label{ex:L08-SS}
Three groups: group 1 $=\{5,7\}$, group 2 $=\{8,10,12\}$, group 3 $=\{4,6\}$. So $k=3$, $n=7$.
\begin{itemize}
  \item Group means $\bar Y_{1\cdot}=6$, $\bar Y_{2\cdot}=10$, $\bar Y_{3\cdot}=5$. Grand mean $\bar Y_{\cdot\cdot}=52/7=7.4286$.
  \item $\SSW=[(5-6)^2+(7-6)^2]+[(8-10)^2+0+(12-10)^2]+[(4-5)^2+(6-5)^2]=2+8+2=12$.
  \item $\bar Y_{i\cdot}-\bar Y_{\cdot\cdot}=-10/7,\ 18/7,\ -17/7$, so
  $\SSB=\frac{2\cdot100+3\cdot324+2\cdot289}{49}=\frac{1750}{49}=\frac{250}{7}=35.714$.
  \item $\TSS=\SSW+\SSB=12+250/7=334/7=47.714$. Directly: $\sum Y_{ij}^2-n\bar Y_{\cdot\cdot}^2=434-2704/7=334/7$. \checkmark
  \item Degrees of freedom $n-k=4$ and $k-1=2$. The ratio $\frac{\SSB/2}{\SSW/4}=\frac{125/7}{3}=\frac{125}{21}=5.952$ is the $F$ statistic of Week~8. Under equal means it has the $F_{2,4}$ distribution. Here the p-value is $0.063$.
\end{itemize}
\end{example}

\subsection{The model in matrix form}
Stack the data group by group:
$Y=(Y_{11},\dots,Y_{1n_1},\,Y_{21},\dots,Y_{2n_2},\,\dots,\,Y_{k1},\dots,Y_{kn_k})\T\in\RR^n$,
and similarly $\eps$. The parameter vector is $\beta=(\mu,\mu_1,\dots,\mu_k)\T\in\RR^{k+1}$ and
the \vocab{design matrix} is the $n\times(k+1)$ matrix
\[
X=\begin{pmatrix}
\one_{n_1} & \one_{n_1} & 0 & \cdots & 0\\
\one_{n_2} & 0 & \one_{n_2} & \cdots & 0\\
\vdots & & & \ddots & \\
\one_{n_k} & 0 & 0 & \cdots & \one_{n_k}
\end{pmatrix}.
\]
Then $Y_{ij}=\mu+\mu_i+\eps_{ij}$ for all $i,j$ is the same as $Y=X\beta+\eps$, and $\EE[Y]=X\beta$.

\begin{proposition}\label{prop:L08-rank}
If every $n_i\ge 1$, then $\rank(X)=k$, while $X$ has $k+1$ columns.
\end{proposition}
\begin{proof}
Let $c_0$ be the first column and $c_1,\dots,c_k$ the indicator columns of the groups. Every row
has exactly one $1$ among $c_1,\dots,c_k$, so $c_0=c_1+\dots+c_k$ and $\rank(X)=\rank[c_1,\dots,c_k]$.
The columns $c_1,\dots,c_k$ are non-zero ($n_i\ge1$) and have disjoint supports. If
$\sum a_ic_i=0$, then looking at a row of group $i$ gives $a_i=0$. So they are linearly
independent, and $\rank(X)=k$.
\end{proof}

\subsection{Identifiability}
How many parameters are there? $k+1$ (plus $\sigma^2$). Can we estimate all of them? The notes'
answer: \emph{no!} Fix any $c\in\RR$ and set
\[
\tilde\mu=\mu+c,\qquad \tilde\mu_i=\mu_i-c\quad(1\le i\le k).
\]
Then $\tilde\mu+\tilde\mu_i+\eps_{ij}=\mu+\mu_i+\eps_{ij}=Y_{ij}$. The two parameter vectors give
\emph{the same} distribution of the data: ``the models are the same!'' No procedure based on the
data can tell $\mu$ from $\mu+c$.

\begin{definition}[Identifiability]\label{def:L08-ident}
A parameter, or a function $\psi(\beta)$ of the parameters, is \vocab{identifiable} if different
values of it always produce different distributions of the data. In the linear model with
$\EE[Y]=X\beta$, this means: $X\beta_1=X\beta_2\Rightarrow\psi(\beta_1)=\psi(\beta_2)$.
\end{definition}

\begin{example}[What is identifiable in Model 1]\label{ex:L08-ident}
\begin{itemize}
  \item $\mu$ is not identifiable: $\beta_1=(\mu,\mu_1,\dots)$ and $\beta_2=(\mu+c,\mu_1-c,\dots)$ have $X\beta_1=X\beta_2$.
  \item $\mu+\mu_i$ is identifiable. It is the mean of group $i$, which is an entry of $X\beta$.
  \item $\mu_1-\mu_2$ is identifiable. It equals $(\mu+\mu_1)-(\mu+\mu_2)$, a difference of entries of $X\beta$. Moreover $Y_{11}-Y_{21}$ is an \emph{unbiased} estimate: $\EE[Y_{11}-Y_{21}]=(\mu+\mu_1)-(\mu+\mu_2)=\mu_1-\mu_2$.
\end{itemize}
The notes ask: is $Y_{11}-Y_{21}$ \emph{optimal}? What is an optimal estimate? This is answered in
Week~5--6 (BLUEs, the Gauss--Markov theorem). The next note gives a first comparison.
\end{example}

\begin{note}[Supplementary: a better unbiased estimate of $\mu_1-\mu_2$]
$\bar Y_{1\cdot}-\bar Y_{2\cdot}$ is also unbiased for $\mu_1-\mu_2$, and
$\Var(\bar Y_{1\cdot}-\bar Y_{2\cdot})=\sigma^2(1/n_1+1/n_2)\le 2\sigma^2=\Var(Y_{11}-Y_{21})$.
The inequality is strict as soon as $n_1$ or $n_2$ exceeds $1$. Using all the data in each group
beats using one observation. The Gauss--Markov theorem shows that $\bar Y_{1\cdot}-\bar Y_{2\cdot}$ is
in fact the best linear unbiased estimate.
\end{note}

The second way out is a \emph{constraint}. Suppose we assume $\sum_{i=1}^k\mu_i=0$. Then all
parameters become identifiable (``Yay!'' in the notes).

\begin{proposition}\label{prop:L08-constraint}
Under the constraint $\sum_{i=1}^k\mu_i=0$, the map $\beta\mapsto X\beta$ is one-to-one on
$\{\beta:\sum\mu_i=0\}$. Explicitly, if $m_i=\mu+\mu_i$ are the group means, then
$\mu=\frac1k\sum_i m_i$ and $\mu_i=m_i-\frac1k\sum_{l}m_l$.
\end{proposition}
\begin{proof}
$X\beta$ determines the group means $m_i=\mu+\mu_i$. Summing, $\sum_i m_i=k\mu+\sum\mu_i=k\mu$, so
$\mu=\frac1k\sum m_i$, and then $\mu_i=m_i-\mu$. Thus $\beta$ is a function of $X\beta$, and any two
constrained $\beta$ with the same $X\beta$ are equal.
\end{proof}

\begin{note}
In either approach, only $k$ parameters (or combinations of parameters) are identifiable:
the $k$ group means $\mu+\mu_i$ and everything computable from them. The notes ask
\emph{why $k$?} Because $\rank(X)=k$ (\cref{prop:L08-rank}). $X\beta$ ranges over a
$k$-dimensional space, so the data can pin down at most $k$ independent linear combinations of
$\beta$. \Cref{sec:L10} makes this precise.
\end{note}

\subsection{Computation}
The notebook \nb{week03/L08\_oneway\_sums\_of\_squares.ipynb}:
\begin{itemize}
  \item recomputes \cref{ex:L08-SS} and checks it against \texttt{statsmodels}' \texttt{anova\_lm};
  \item builds $X$ for Model~1 and checks $\rank(X)=k$;
  \item shows non-identifiability numerically ($X\beta_1=X\beta_2$ for $\beta_2=\beta_1+c(1,-1,\dots,-1)$);
  \item simulates Model~1 many times and compares the histograms of $\SSW/\sigma^2$ and $\SSB/\sigma^2$ with the $\chi^2_{n-k}$ and $\chi^2_{k-1}$ densities, and checks that their correlation is near $0$.
\end{itemize}

\subsection{Summary}
\begin{itemize}
  \item Model 1: $Y_{ij}=\mu+\mu_i+\eps_{ij}$, $\eps_{ij}$ i.i.d.\ $\N(0,\sigma^2)$, $n=\sum n_i$.
  \item $\TSS=\SSW+\SSB$, with degrees of freedom $n-1=(n-k)+(k-1)$.
  \item $\SSW/\sigma^2\sim\chi^2_{n-k}$ and $\SSW\indep\SSB$. Under equal means, $\SSB/\sigma^2\sim\chi^2_{k-1}$.
  \item $Y=X\beta+\eps$ with $\beta\in\RR^{k+1}$ and $\rank(X)=k$. The parameters are not identifiable ($\mu\to\mu+c$, $\mu_i\to\mu_i-c$).
  \item Group means $\mu+\mu_i$ and contrasts $\mu_i-\mu_j$ are identifiable. The constraint $\sum\mu_i=0$ makes everything identifiable.
\end{itemize}

\subsection{Exercises}
\begin{problem}\label{prob:L08-1}
Show that $S_i^2$ is an unbiased estimate of $\sigma^2$, and that $\SSW/(n-k)$ is unbiased for $\sigma^2$ (without assuming normality).
\end{problem}
\begin{problem}\label{prob:L08-2}
For the groups $\{1,3\}$ and $\{2,4,6\}$ compute $\SSW$, $\SSB$ and $\TSS$ and check \cref{thm:L08-decomp}.
\end{problem}
\begin{problem}\label{prob:L08-3}
Show that in Model 1 (normality not needed) $\EE[\SSB]=(k-1)\sigma^2+\sum_i n_i(\mu_i-\bar\mu_w)^2$, where $\bar\mu_w=\frac1n\sum_i n_i\mu_i$.
\end{problem}
\begin{problem}\label{prob:L08-4}
Prove that the trace of a symmetric idempotent matrix equals its rank.
\end{problem}
\begin{problem}\label{prob:L08-5}
Under the constraint $\sum_i n_i\mu_i=0$ (instead of $\sum_i\mu_i=0$), express $\mu$ and $\mu_i$ in terms of the group means.
\end{problem}

\subsection{Further reading}
Christensen, \S1.2--1.3 (multivariate normal, distribution of quadratic forms), \S2.1
(identifiability and estimability) and Chapter~4 (one-way ANOVA). Rao, \S3b (quadratic forms
in normal variables) and \S4a.
